\documentclass[11pt,a4paper]{article}
% Dynamics note (P0 of the research plan "Bits of Evidence", version 7.2).
% Alignment Forum version. Compile with pdflatex (run twice for references).
% Every number in this note is taken from the plan's Appendix A; the scripts
% that produced them are archived (see the Reproducibility paragraph).
\usepackage[utf8]{inputenc}
\usepackage[T1]{fontenc}
\usepackage[margin=2.5cm]{geometry}
\usepackage{amsmath,amssymb,booktabs,xcolor}
\usepackage[colorlinks=true,linkcolor=blue!55!black,citecolor=blue!55!black,urlcolor=blue!55!black,pdftitle={Learning coefficients predict selection, not passage},pdfauthor={Sven Heydenreich}]{hyperref}

% --- notation -------------------------------------------------------------
\newcommand{\E}{\mathbb E}
\newcommand{\Var}{\operatorname{Var}}
\newcommand{\eps}{\varepsilon}
\newcommand{\dd}{\mathrm d}
\newcommand{\chiexit}{\chi_{\mathrm{exit}}}
\newcommand{\unv}{\textsuperscript{\dag}}

% --- evidence labels, as in the research plan -----------------------------
\definecolor{cknown}{RGB}{20,100,40}
\definecolor{ccalc}{RGB}{25,70,150}
\definecolor{cnum}{RGB}{0,105,115}
\definecolor{chyp}{RGB}{150,30,40}
\newcommand{\evl}[2]{{\small\bfseries\color{#1}[#2]}}
\newcommand{\LKnown}{\evl{cknown}{Known result}}
\newcommand{\LCalc}{\evl{ccalc}{Internal calculation}}
\newcommand{\LNum}{\evl{cnum}{Numerical check}}
\newcommand{\LHyp}{\evl{chyp}{Hypothesis}}

\setlength{\emergencystretch}{2em}
\setlength{\parskip}{0.25em}

\title{\textbf{Learning coefficients predict selection, not passage}\\[0.4em]
\large What the singular-learning-theory invariants of a flat set do and do not say\\
about how long stochastic gradient descent stays there}
\author{Sven Heydenreich}
\date{October 2026}

\begin{document}
\maketitle

\begin{abstract}
\noindent
Singular learning theory attaches to every flat set of the population loss
two invariants, the learning coefficient $\lambda$ and its multiplicity $m$,
and these invariants fix the posterior mass and the Bayesian code length of
the set. It is tempting to read them as dynamical quantities as well: the
flatter a saddle in the sense of $\lambda$, the longer SGD should linger
there. This note argues that the reading is wrong as a general statement and
that it fails in checkable ways in the simplest models where it can be
tested. Four controls hold $(\lambda,m)$ fixed and change exit times by
changing the teacher, the gradient covariance, the initialization scales or
the metric. For isotropic diffusion through a flat set, the published
passage law $T\propto\eps^\lambda$ is replaced by
$T\propto\eps^{\lambda-\chiexit}$ up to logarithms, where $\chiexit$ is the
scaling exponent of the exit cross-section; three worked germs give
constant, logarithmic and $\eps^{-1/4}$ passage times where $\eps^\lambda$
predicts decay, and finite-difference solutions agree. For a discrete
balanced chain, the diffusion approximation is accurate only for small
steps, its strong-noise power law never appears, label noise traps about
$0.9\%$ of paths forever, and exit time is not a function of the ratio of
step size to batch size. The positive lesson is that $\lambda$ governs
selection and evidence, that is, which region the posterior prefers and how
many bits it costs, and not the time it takes a training process to pass.
\end{abstract}

\paragraph{Evidence labels.} Following the research plan this note is drawn
from, claims are tagged. \LKnown{} is established in the literature and
cited where used. \LCalc{} is worked out here or in the companion appendix
under stated assumptions; the label implies neither novelty nor a completed
proof. \LNum{} is a completed computation on the stated mathematical model,
archived with code; it says nothing about observed training of large
networks. \LHyp{} is a claim the author expects to be true and has not
established.

\section{Summary in plain language}
\label{sec:summary}

Singular learning theory (SLT) says that the Bayesian posterior at sample
size $n$ gives a flat set $S$ of the population loss a weight governed by
its learning coefficient $\lambda_S$: the free energy of a neighbourhood of
$S$ is $nL_n+\lambda_S\log n-(m_S-1)\log\log n+O(1)$ \cite{watanabe2009}.
Smaller $\lambda_S$ means more posterior mass and a shorter code. In
developmental interpretability it has become common to also read
$\lambda_S$ dynamically, as a measure of how long gradient-based training
is held up at $S$. A recent paper makes this precise for diffusion through
a flat set and derives a passage time proportional to $\eps^{\lambda}$,
where $\eps$ is the loss level that defines the set \cite{hennick2026}.

This note makes three claims and one observation.

\begin{enumerate}
\item \LCalc{} \emph{Scoped impossibility.} No function of the saddle
invariants $(\lambda_s,m_s)$, and indeed no statistic of the population loss
alone, determines the statistics of saddle passage once the teacher, the
gradient covariance, the initialization scales or the metric are allowed to
vary. Four explicit controls keep the invariants fixed and change the exit
time by polynomial factors (Section~\ref{sec:impossibility}).
\item \LCalc{} and \LNum{} \emph{Passage through flat sets.} For isotropic
diffusion in a tube $\{K<\eps\}$ with reflecting walls, the Fick--Jacobs
reduction gives $T\propto\eps^{\lambda-\chiexit}$ up to logarithms, where
$\chiexit$ is the scaling exponent of the exit cross-section. In three
germs with $\lambda\in\{\tfrac12,\tfrac12,\tfrac14\}$ the passage time is
constant, logarithmic in $1/\eps$, and proportional to $\eps^{-1/4}$; the
law $\eps^\lambda$ predicts decay in all three. Finite-difference solutions
of the full two-dimensional problem agree (Section~\ref{sec:flat}).
\item \LNum{} \emph{Discrete chain.} For the discrete balanced update, the
continuous approximation is accurate for small steps and its strong-noise
power law never appears. With label noise about $0.9\%$ of paths cross zero
and never exit, so the unconditional mean exit time is infinite. Exit time
is not a function of $\eta/b$ alone (Section~\ref{sec:discrete}).
\item \emph{The positive lesson.} $\lambda$ governs posterior masses and code
lengths, that is, selection and evidence, not passage times. A
basin-selection account of grokking makes the same distinction in a
different setting \cite{cullen2026} (Section~\ref{sec:lesson}).
\end{enumerate}

What the note does not claim. It says nothing about training of large
networks; every computation concerns a stated low-dimensional model. It does
not say that $\lambda$ is irrelevant to dynamics: under a fixed optimizer
with a fixed teacher, finite-temperature quantities such as activated escape
rates do depend on $\lambda$ through the equilibrium masses of wells. It
does not claim that the published law is wrong in the authors' exact
formulation; the calculation here uses a reflecting-tube caricature, and the
discrepancy will be checked with the authors, who receive a courtesy copy of
this note before it is posted.

\section{The question}
\label{sec:question}

Let $K(w)\ge0$ be the population loss of a model with parameter $w$, and let
$S$ be a connected component of its zero set or, more generally, of a level
set of critical points that the training trajectory has to leave. \LKnown{}
Near $S$ the volume of the sublevel set scales as
\begin{equation}
\operatorname{vol}\{w\ \text{near}\ S:\ K(w)<\eps\}\;\asymp\;
\eps^{\lambda_S}\,\bigl(\log\tfrac1\eps\bigr)^{m_S-1},\qquad \eps\to0,
\label{eq:volume}
\end{equation}
and the same pair $(\lambda_S,m_S)$ appears in the local free energy
\begin{equation}
F_U(n)=-\log\int_U e^{-nL_n(w)}\varphi(w)\,\dd w
= nL_n(w_0)+\lambda_U\log n-(m_U-1)\log\log n+c_U+o(1)
\label{eq:localF}
\end{equation}
of any neighbourhood $U$ of $S$ \cite{watanabe2009}. Here $L_n$ is the
empirical average negative log likelihood, $w_0\in S$ an optimal parameter
and $\varphi$ a prior with positive density near $S$. The posterior odds
between two regions are $\exp[-(F_{U_1}-F_{U_2})]$, and the expected excess
log loss of the Bayes predictive at sample size $n$ is $\lambda/n$. Local
values of $\lambda$ at a trained parameter are estimated by stochastic
gradient Langevin dynamics \cite{lau2025llc}.

The invariants are therefore volume exponents. Training by SGD is a
process, and the question is whether these volume exponents also control
the process. Two forms of the question are in circulation.

\begin{description}
\item[Saddle passage.] A trajectory approaches a saddle set $S$ with
invariants $(\lambda_s,m_s)$ and leaves after a plateau. Does some function
of $(\lambda_s,m_s)$ determine the distribution of the plateau length?
\item[Passage through a flat set.] A diffusion (the continuous caricature of
noisy SGD) moves inside $\{K<\eps\}$ until it reaches an escape set. Does
the mean passage time scale as $\eps^{\lambda}$ \cite{hennick2026}?
\end{description}

Both questions are answered in the negative below, in the strongest form
the author could make checkable. The negative answer is not a technicality:
the quantities that control the process are the ones the volume exponent
discards, namely the direction and strength of the drift, the covariance of
the noise, the sizes of the cross-sections through which the trajectory
must pass, and the step size.

\section{Scoped impossibility: four controls}
\label{sec:impossibility}

\LCalc{} The controls use the balanced linear chain of depth $L$: a product
of $L$ scalar weights that stay equal under gradient descent, so the state is
one number $u$ and the student output is $u^Lx$. The population loss has a
flat set at $u=0$ whose invariants are $(\lambda_s,m_s)=(1,L)$ for a
realizable teacher. Each control holds the model-intrinsic invariants fixed
and changes one ingredient of the training process; in each case the exit
time changes by a polynomial factor or more. The derivations were given in
version~5 of the author's research plan (its Appendices A and B and
Section~6.2) and are not reproduced here; they will appear in a companion
appendix to the arXiv version of this note.

\begin{center}
\small
\begin{tabular}{@{}p{2.1cm}p{3.6cm}p{3.5cm}p{4.9cm}@{}}
\toprule
Control & Held fixed & Changed & Consequence for exit \\
\midrule
Teacher & Self-divergence and $(\lambda_s,m_s)=(1,L)$ & Realizable
teacher, cancelling versus non-cancelling & Exit time
$\propto u_0^{-(2L-2)}$ versus $u_0^{-(L-2)}$ \\ \addlinespace
Gradient covariance & Population loss, Fisher, self-divergence &
Gaussian versus Rademacher inputs & Positive versus exactly zero
minibatch noise \\ \addlinespace
Scale hierarchy & Self-divergence, teacher & Initialization
$(\eps,\eps,\sqrt\eps,1)$ versus equal small layers &
$\tfrac12\eps^{-1/2}\log(1/\eps)$ versus $\eps^{-(r-2)}$ \\ \addlinespace
Metric & Population loss, self-divergence & Euclidean versus natural
gradient & Polynomial delay versus bounded time; fixed damping restores
the delay \\
\bottomrule
\end{tabular}
\end{center}

\noindent Here $u_0$ is the initial scale of the chain.\footnote{The
``self-divergence'' and the exponent $r$ in the scale-hierarchy row are
defined in the version-5 derivations. \textbf{[TODO: import the definitions
of the self-divergence and of $r$ from the version-5 appendix before
posting.]}}

\paragraph{Reading the table.} The teacher control is the sharpest. Two
realizable teachers give the same population loss near the saddle up to the
invariants, but one of them makes the leading drift term cancel and the
other does not, and the exit time changes from $u_0^{-(L-2)}$ to
$u_0^{-(2L-2)}$. The gradient-covariance control keeps the population loss,
the Fisher information and the self-divergence fixed and switches the input
distribution from Gaussian to Rademacher: with Rademacher inputs the
minibatch gradient of the chain has exactly zero variance, so the
``noise-driven'' escape of the Gaussian case does not exist. The
scale-hierarchy control shows that the initialization, not the loss,
decides between a logarithmically corrected $\eps^{-1/2}$ delay and a
higher power. The metric control shows that natural-gradient descent on the
same loss leaves the saddle in bounded time, and that adding a fixed damping
to the natural metric brings the polynomial delay back.

\paragraph{What the impossibility covers and what it does not.} The
statement is: no function of $(\lambda_s,m_s)$, nor of any other statistic
of the population loss alone, determines saddle-passage statistics once
teacher, gradient covariance, initialization scales or metric may vary.
This is a statement about the class of training processes that share a
population loss. It does \emph{not} cover teacher-dependent
finite-temperature quantities under a fixed optimizer. In particular, for an
activated escape from a well at temperature $T$ under a Gibbs law, the
escape rate is inversely proportional to the equilibrium mass of the well,
and by \eqref{eq:volume} that mass scales as $T^{\lambda}$ up to
logarithms; so $\lambda$ does enter activated escape rates, through the
masses of wells, which is exactly the role it plays in selection. The claim
of this note is that $\lambda$ enters the dynamics \emph{only} in that
role, not that it never appears.

\section{Passage through flat sets}
\label{sec:flat}

\LKnown{} Hennick and De Baerdemacker argue that the time for isotropic
diffusion to pass through a small flat region scales as $\eps^\lambda$
divided by the capacitance of the escape set, which is treated as
independent of $\eps$ \cite{hennick2026}. The argument has the attractive
feature that it ties a training-time quantity to the volume exponent of
\eqref{eq:volume}: the diffusion must ``fill'' a region of volume
$\eps^\lambda$ before it finds the exit.

\LCalc{} Consider diffusion with coefficient $D$ in the tube
$W_\eps=\{K<\eps\}$ inside a fixed domain, reflecting on $\{K=\eps\}$ and
absorbed at an escape set. For thin tubes the Fick--Jacobs reduction
\cite{zwanzig1992,reguera2001} gives a one-dimensional problem along the
zero set, with cross-section $A(s)$ and enclosed volume
$V(s)=\int_0^sA(u)\,\dd u$. Starting at the centre, with reflection there by
symmetry, the mean first-passage time is
\begin{equation}
T\approx\frac1D\int_0^{s_{\rm exit}}\frac{V(s)}{A(s)}\,\dd s .
\label{eq:fj}
\end{equation}
The derivation is in Appendix~\ref{app:fj}. The reduction is the
standard tool for diffusion through narrow channels with an entropic
barrier: the density equilibrates quickly across each cross-section, and
the slow motion along the channel feels the cross-section only through the
``entropic potential'' $-\log A(s)$.

Formula \eqref{eq:fj} already contains the answer. The volume $V$ is
dominated by the bulk of the tube and scales as $\eps^\lambda$ up to
logarithms, in agreement with \eqref{eq:volume}. But the integrand divides
by the cross-section, and near the exit the cross-section is the exit
window itself, which scales as $\eps^{\chiexit}$ for some exponent
$\chiexit$ that depends on how the zero set meets the escape set. The
integral is dominated near the exit, giving
\begin{equation}
T\;\propto\;\eps^{\lambda-\chiexit}\qquad\text{up to logarithms.}
\label{eq:scaling}
\end{equation}
The capacitance of an escape window that is itself a cross-section of the
tube scales with $\eps$; it is not constant. Three germs make this
concrete, with $D=1$ for the numbers (details in
Appendix~\ref{app:germs}).
\begin{enumerate}
\item $K=y^2$ in the unit disk ($\lambda=\tfrac12$). Here $A=2\sqrt\eps$
and $V=2\sqrt\eps\,x$, so $T=\int_0^1x\,\dd x=1/(2D)$, independent of
$\eps$.
\item $K=x^2y^2$ in the unit disk ($\lambda=\tfrac12$, $m=2$). Here
$A(x)=2\min(x,\sqrt\eps/x)$, $V(x)=\sqrt\eps+2\sqrt\eps\log(x\eps^{-1/4})$
for $x>\eps^{1/4}$, and $T=\log(1/\eps)/(8D)+O(\sqrt\eps)$. The
multiplicity survives as a logarithm.
\item $K=x^4y^2$ in $[-1,1]^2$, absorbed only at $x=\pm1$ ($\lambda=\tfrac14$).
Here $A(x)=\min(2,2\sqrt\eps/x^2)$, $V(x)=4\eps^{1/4}-2\sqrt\eps/x$ for
$x>\eps^{1/4}$, and $T=\tfrac23\eps^{-1/4}/D-1/(2D)+O(\sqrt\eps)$: the
time grows as $\eps\to0$.
\end{enumerate}
All three cases have $\chiexit=\tfrac12$, because in each the exit window
is a cross-section of a tube of width $\sqrt\eps$ around a smooth branch of
the zero set. The exponents $\lambda-\chiexit$ are $0$, $0$ and $-\tfrac14$,
and the law $\eps^\lambda$ predicts a decreasing passage time in all three.

\LNum{} Finite-difference solutions of the reflected-diffusion
mean-first-passage problem in the full two-dimensional tube
(\texttt{tube\_mfpt.py}) agree with \eqref{eq:fj}. The columns ``reduction''
list the closed forms above without their $O(\sqrt\eps)$ remainders.

\begin{center}\small
\begin{tabular}{@{}ccccccc@{}}
\toprule
$\eps$ & (1) numerical & (2) numerical & (2) reduction & (3) numerical & (3) reduction & $\eps^\lambda$ law \\
\midrule
$10^{-2}$ & 0.500 & 0.620 & 0.576 & 1.835 & 1.608 & decreasing \\
$10^{-3}$ & 0.504 & 0.883 & 0.863 & 3.434 & 3.249 & decreasing \\
$10^{-4}$ & 0.500 & 1.150 & 1.151 & 6.253 & 6.167 & decreasing \\
$10^{-5}$ & --- & 1.430 & 1.439 & 11.359 & 11.355 & decreasing \\
\bottomrule
\end{tabular}
\end{center}

\noindent The value 0.504 reflects grid alignment. In case (3) the fitted
log--log slope is $-0.264$ overall and $-0.259$ on the last decade, against
the predicted $-1/4$; in case (2) the slope of $T$ against $\log(1/\eps)$
is $0.117$, against $1/8$. The agreement improves as $\eps$ decreases, as
it should for an asymptotic reduction.

\paragraph{Caveat.} The calculation uses a reflecting-tube caricature:
isotropic diffusion confined to $\{K<\eps\}$ by a hard wall. The exact
formulation of \cite{hennick2026} differs in how the region is defined and
how the diffusion is confined, and whether the correction transfers to it
will be checked with the authors, who receive a courtesy copy before this
note is posted. The result is consistent with the order-one motion along
zero sets found for Langevin dynamics of real-analytic potentials
\cite{larsson2026}: motion along a flat set is not slow in the sense that
the $\eps^\lambda$ law requires.

\section{The discrete balanced chain}
\label{sec:discrete}

\LNum{} The continuous caricature is itself a step away from SGD, which
takes discrete steps of size $\eta$ on minibatches of size $b$. The
balanced update for the depth-$L$ chain with teacher $sx+q$ is
\begin{equation}
u_{k+1}=u_k+\eta\,u_k^{L-1}\bigl[(s-u_k^L)A_k+C_k\bigr],\qquad
A_k=\frac1b\sum_{i=1}^b x_i^2,\quad C_k=\frac1b\sum_{i=1}^b q_ix_i,
\label{eq:update}
\end{equation}
derived in Appendix~\ref{app:discrete}. The checks use $L=3$, $s=1$,
Gaussian inputs, the initial value $w_0=10^{-3}$, the exit level
$w_1=1/2$, and the clock $T=K\eta$ with $K$ the number of steps
(\texttt{discrete\_sgd\_check.py}; 4000 paths per setting). The row
``diffusion'' is the mean exit time of the continuous approximation,
computed by exact quadrature; the row ``independent re-run'' is from an
external reviewer of an earlier version of the research plan.

\begin{center}\small
\begin{tabular}{@{}lcccccc@{}}
\toprule
$\eta$ (with $b=1$, $\sigma_q=0$) & 0.01 & 0.1 & 1 & 3 & 10 & 100 \\
\midrule
This check & 9.19 & 9.75 & 13.98 & 20.70 & 38.3 & 179.0 \\
Standard error & 0.007 & 0.021 & 0.063 & 0.12 & 0.24 & 1.3 \\
Independent re-run & 9.19 & 9.79 & 14.05 & 20.7 & 38.7 & 182 \\
Diffusion, exact quadrature & 9.15 & 9.51 & 14.25 & --- & 186.9 & 14278 \\
\bottomrule
\end{tabular}
\end{center}

\noindent Three things follow. First, without label noise $A_k\ge0$, so
every step moves $u$ towards the teacher and no path ever moves backward:
the ``noise'' in this regime only modulates the speed. Second, the
diffusion is accurate up to $r=\eta/b\approx1$ and then fails
qualitatively. Its strong-noise power law, visible in the last two columns
of the quadrature row, never appears in discrete SGD: at $\eta=100$ the
discrete chain exits in fewer than two steps on average ($T/\eta\approx1.8$),
so $T$ grows only linearly with $\eta$. Third, the small-step limit is
reached smoothly from above, so the continuous approximation is a good
description of the chain exactly where it is least interesting.

\paragraph{Label noise and trapping.} \LNum{} With label noise
($\sigma_q=1$, $\eta=1$, $b=1$; \texttt{trapped\_check.py}; 20{,}000 paths,
$10^6$ steps), 179 paths did not exit ($0.90\%\pm0.07\%$) and 212 crossed
zero. Every non-exiting path sits near
$u\approx-1.0\times10^{-6}\approx-1/(\eta sk)$ after $k=10^6$ steps. The
mechanism is visible in \eqref{eq:update} with $L=3$: for $u<0$ the drift
$\eta su^2$ pushes $u$ towards zero from below ever more slowly, since the
drift vanishes quadratically at the flat set, and returning above zero
would require a noise kick $\eta u^2C_k$ of size $|u|$, that is, $|C_k|$ of
order $1/(\eta|u|)$, which becomes astronomically unlikely as $u\to0^-$.
The median exit time is 15 and the mean conditional on exit about 16.3;
the unconditional mean is infinite. A reader who summarizes this
distribution by its mean reports nothing; a reader who summarizes it by a
function of $(\lambda_s,m_s)$, which are the same with and without label
noise, misses the trapped population entirely.

\paragraph{Exit time is not a function of $\eta/b$.} \LNum{} With
$r=\eta/b=0.1$ fixed, $(\eta,b)=(0.1,1)$, $(1,10)$ and $(10,100)$ give
$T=9.77$, $11.61$ and $29.71$ (independent re-run: $11.8$ and $29.6$ for
the last two). The continuous approximation depends on $\eta$ and $b$ only
through $r$; the discrete chain does not.

\section{The positive lesson}
\label{sec:lesson}

What survives is the role that SLT assigns to $\lambda$ in the first place.
\LKnown{} By \eqref{eq:localF}, $\lambda_U$ and $m_U$ fix the posterior mass
of a region up to a prior-dependent constant, so they decide which of two
regions the posterior prefers at sample size $n$, and by how much. The same
quantities fix the redundancy of the Bayes mixture code,
$F_n-nL_n(w_0)=\lambda\log n-(m-1)\log\log n+O(1)$, and the expected excess
log loss $\lambda/n$ of the Bayes predictive \cite{watanabe2009}. These are
statements about selection and evidence: about which solution the data
support and how many bits that support is worth.

Passage is a different kind of quantity. It depends on the drift along the
flat set, on the noise covariance, on the cross-sections the process has to
pass through, and on the step size, none of which are functions of the
population loss alone, and all of which the four controls of
Section~\ref{sec:impossibility} vary at fixed $(\lambda,m)$. Even the
stationary law of SGD near a minimum is not the Gibbs law that SLT
describes: for gradient noise shaped like the curvature the stationary
covariance solves $H\Sigma+\Sigma H=(\eta/B)H$ \cite{mandt2017}, and a
loss-fluctuation estimator of $\lambda$ applied to such a chain returns a
curvature statistic rather than the learning coefficient. Estimating
$\lambda$ from a sampler is therefore a question about the sampler
\cite{lau2025llc,hitchcock2025,cai2026}, and reading $\lambda$ off a
training trajectory is not the same as measuring it.

\LHyp{} The author expects the same distinction to organize empirical
results in developmental interpretability: where the local learning
coefficient tracks a phase transition, it does so because the transition is
a change in which region the posterior (or the tempered local posterior
that the estimator samples) prefers, and not because $\lambda$ sets the
time at which the optimizer crosses. A basin-selection account of grokking
\cite{cullen2026} makes the same distinction in a different setting: the
delayed generalization is explained by which basin the posterior selects at
a given sample size, not by the passage time to it. That account and this
note are complementary; neither predicts how long the crossing takes.

\paragraph{Reproducibility.} All numbers in this note are taken from the
research plan the note is drawn from. The numerical checks were produced by
the archived scripts named there (\texttt{tube\_mfpt.py},
\texttt{discrete\_sgd\_check.py}, \texttt{trapped\_check.py}) and will be
released with the note. The four controls of Section~\ref{sec:impossibility}
are stated with their results only; their derivations will be added as a
companion appendix to the arXiv version. The research plan marks
bibliographic entries not independently verified with a dagger\unv{}; none
of the entries cited here carries that mark.

\appendix
\renewcommand{\theequation}{A.\arabic{equation}}
\setcounter{equation}{0}

\section{Derivations}
\label{app:derivations}

\subsection{The Fick--Jacobs reduction and the passage-time integral}
\label{app:fj}

\LKnown{} Let a particle diffuse with coefficient $D$ in a thin tube whose
axis is parametrized by $s$ and whose cross-section at $s$ has measure
$A(s)$, with reflecting walls. If the tube is thin enough that the density
equilibrates across each cross-section much faster than it moves along the
axis, the marginal density $p(s,t)$ along the axis obeys the Fick--Jacobs
equation
\begin{equation}
\partial_tp=D\,\partial_s\Bigl[A(s)\,\partial_s\frac{p}{A(s)}\Bigr],
\label{eq:fjpde}
\end{equation}
that is, one-dimensional diffusion in the entropic potential
$U(s)=-\log A(s)$ \cite{zwanzig1992}; corrections for walls that are not
slowly varying are discussed in \cite{reguera2001}. The generator is
$\mathcal L=D\,e^{U}\partial_s(e^{-U}\partial_s)$ and the mean first-passage
time $T(s_0)$ from $s_0$ to an absorbing point $s_{\rm exit}$, with a
reflecting point at $s=0$, solves $\mathcal LT=-1$ with $T'(0)=0$ and
$T(s_{\rm exit})=0$. Integrating twice,
\begin{equation}
T(s_0)=\frac1D\int_{s_0}^{s_{\rm exit}}\frac{\dd s}{A(s)}\int_0^{s}A(u)\,\dd u
=\frac1D\int_{s_0}^{s_{\rm exit}}\frac{V(s)}{A(s)}\,\dd s,
\qquad V(s)=\int_0^sA(u)\,\dd u.
\label{eq:mfpt1d}
\end{equation}
\LCalc{} In the flat-set application the zero set extends to both sides of
the starting point, and a start at the centre is equivalent, by symmetry, to
a reflecting boundary at $s=0$; setting $s_0=0$ gives \eqref{eq:fj}. The
reduction treats $\{K<\eps\}$ as a tube around a branch of the zero set
whose cross-section at $s$ is the set of transverse displacements with
$K<\eps$; for a germ $K=f(s)\,y^2$ this is the interval
$|y|<\sqrt{\eps/f(s)}$, so $A(s)=2\sqrt{\eps/f(s)}$ wherever this is
smaller than the transverse extent of the domain.

The scaling \eqref{eq:scaling} follows from \eqref{eq:mfpt1d}. The total
volume $V(s_{\rm exit})$ is the volume of the tube, which by
\eqref{eq:volume} scales as $\eps^\lambda(\log1/\eps)^{m-1}$. Near the
exit, $V(s)$ is of this order while $A(s)$ is of the order of the exit
window, $\eps^{\chiexit}$. If $V/A$ is integrable up to the exit over a
distance of order one, as in the three germs, the integral is of order
$\eps^{\lambda-\chiexit}$ up to logarithms; if $A$ shrinks faster than $V$
towards the exit, the integral is even more strongly dominated there. In no
case is the dependence on the exit window a constant.

\subsection{The three germs}
\label{app:germs}

\LCalc{} Throughout, $D=1$, $a:=\eps^{1/4}$, and the exit is at $x=1$ along
the positive $x$-axis; the symmetric half-tube contributes the same by
reflection.

\paragraph{Germ 1: $K=y^2$ in the unit disk.} The tube is $|y|<\sqrt\eps$,
so $A(x)=2\sqrt\eps$ for all $x$ up to the disk boundary, to leading order
in $\eps$. Then $V(x)=2\sqrt\eps\,x$, $V/A=x$, and
\begin{equation}
T=\int_0^1x\,\dd x=\tfrac12,
\end{equation}
independent of $\eps$. The volume is $\propto\eps^{1/2}$ and so is the exit
cross-section: $\lambda=\chiexit=\tfrac12$ and the exponent in
\eqref{eq:scaling} vanishes. The numerical column (1) of the table in
Section~\ref{sec:flat} shows $0.500$ at every $\eps$ computed.

\paragraph{Germ 2: $K=x^2y^2$ in the unit disk.} The tube
$\{x^2y^2<\eps\}$ is a cross with two branches along the axes; along the
$x$-branch it is $|y|<\sqrt\eps/|x|$. The branches merge in the central
square $|x|,|y|<a$, which the reduction splits between the two branches
along its diagonals; the part belonging to the $x$-branch at abscissa $x<a$
is $|y|<x$. Hence
\begin{equation}
A(x)=2\min\bigl(x,\sqrt\eps/x\bigr),\qquad
V(x)=\begin{cases}x^2,&x\le a,\\[2pt]
\sqrt\eps+2\sqrt\eps\log(x/a),&x>a,\end{cases}
\end{equation}
using $\int_a^x2\sqrt\eps\,u^{-1}\dd u=2\sqrt\eps\log(x/a)$ and
$a^2=\sqrt\eps$. The central piece of the integral, $\int_0^a(x^2/2x)\,\dd x$,
is $O(\sqrt\eps)$. For $x>a$,
\begin{equation}
\frac{V(x)}{A(x)}=\frac{x}{2}+x\log\frac{x}{a}
=\frac{x}{2}+x\log x+\frac{x}{4}\log\frac1\eps ,
\end{equation}
and integrating from $a$ to $1$: $\int_a^1\tfrac x2\dd x=\tfrac14+O(\sqrt\eps)$,
$\int_a^1x\log x\,\dd x=-\tfrac14+O(\sqrt\eps\log(1/\eps))$, and
$\tfrac14\log(1/\eps)\int_a^1x\,\dd x=\tfrac18\log(1/\eps)+O(\sqrt\eps\log(1/\eps))$.
The constants cancel and
\begin{equation}
T=\tfrac18\log\tfrac1\eps+O(\sqrt\eps).
\end{equation}
The multiplicity $m=2$ of the germ appears as the logarithm in the volume
and survives in $T$; the exit cross-section is again $2\sqrt\eps$, so
$\lambda=\chiexit=\tfrac12$ and the power vanishes.

\paragraph{Germ 3: $K=x^4y^2$ in $[-1,1]^2$, absorbed only at $x=\pm1$.}
The tube along the $x$-axis is $|y|<\sqrt\eps/x^2$, clipped by the box at
$|y|=1$, so
\begin{equation}
A(x)=\min\bigl(2,\,2\sqrt\eps/x^2\bigr),\qquad
V(x)=\begin{cases}2x,&x\le a,\\[2pt]
2a+2\sqrt\eps\,(a^{-1}-x^{-1})=4a-2\sqrt\eps/x,&x>a,\end{cases}
\end{equation}
using $\sqrt\eps/a=a$. The box is reflecting at $y=\pm1$; the only exit is
the cross-section at $x=1$, of width $2\sqrt\eps$. The central piece
$\int_0^a(2x/2)\,\dd x$ is $O(\sqrt\eps)$. For $x>a$,
\begin{equation}
\frac{V(x)}{A(x)}=\frac{(4a-2\sqrt\eps/x)\,x^2}{2\sqrt\eps}
=2\,\eps^{-1/4}x^2-x,
\end{equation}
and $\int_a^1(2\eps^{-1/4}x^2-x)\,\dd x=\tfrac23\eps^{-1/4}-\tfrac12+O(\sqrt\eps)$, so
\begin{equation}
T=\tfrac23\,\eps^{-1/4}-\tfrac12+O(\sqrt\eps).
\end{equation}
Here the volume of the tube is $\propto a=\eps^{1/4}$, in agreement with
$\lambda=\tfrac14$, while the exit cross-section is $\propto\eps^{1/2}$:
$\lambda-\chiexit=-\tfrac14$ and the passage time grows as $\eps\to0$. The
fitted log--log slopes of the two-dimensional solution, $-0.264$ overall
and $-0.259$ on the last decade, approach $-\tfrac14$ from below as the
$-\tfrac12$ offset becomes negligible.

\subsection{The discrete update}
\label{app:discrete}

\LCalc{} Take a linear chain of depth $L$ with scalar weights
$w_1,\dots,w_L$, student output $(\prod_iw_i)x$, teacher output $sx+q$ with
input $x$ and label noise $q$ of standard deviation $\sigma_q$, and the
squared loss $\tfrac12\bigl[(\prod_iw_i)x-sx-q\bigr]^2$ per sample. The
gradient with respect to $w_j$ is
$\bigl(\prod_{i\ne j}w_i\bigr)\,x\,\bigl[(\prod_iw_i)x-sx-q\bigr]$. If all
weights are equal, $w_i=u$, the gradient is the same for every layer, so
equal weights stay equal under gradient descent and the chain is described
by the single number $u$ with student output $u^Lx$. A minibatch step of
size $\eta$ on $b$ samples gives
\begin{equation}
u_{k+1}=u_k-\eta\,u_k^{L-1}\,\frac1b\sum_{i=1}^b x_i\bigl[u_k^Lx_i-sx_i-q_i\bigr]
=u_k+\eta\,u_k^{L-1}\bigl[(s-u_k^L)A_k+C_k\bigr],
\end{equation}
with $A_k=b^{-1}\sum_ix_i^2$ and $C_k=b^{-1}\sum_iq_ix_i$, which is
\eqref{eq:update}. The flat set at $u=0$ has drift of order $u^{L-1}$,
hence the plateau. Its invariants $(\lambda_s,m_s)=(1,L)$ refer to the full
$L$-layer parametrization and are computed in the version-5 derivations;
they are the same with and without label noise, because the label noise
adds only the constant $\tfrac12\sigma_q^2$ to the population loss.
\textbf{[TODO: confirm against the version-5 appendix that $(1,L)$ is the
pair quoted for this saddle in the full parametrization.]}

Two consequences used in Section~\ref{sec:discrete} are immediate. Without
label noise, $C_k=0$ and $A_k\ge0$, so while $0<u_k^L<s$ every step
increases $u_k$ and no path moves backward; the randomness of $A_k$ only
modulates the speed. With label noise and $L=3$, a path that has crossed to
$u<0$ has drift $\eta u^2(s-u^3)A_k\approx\eta su^2A_k>0$ towards zero from
below, which vanishes quadratically, and the noise term $\eta u^2C_k$ must
exceed $|u|$ for the path to return to $u>0$, so the required
$|C_k|\gtrsim1/(\eta|u|)$ grows without bound as the path creeps towards
zero. The observed resting positions $u\approx-1/(\eta sk)$ after $k$ steps
are what the deterministic drift $\dot u=\eta su^2$ gives from any negative
start, since its solution satisfies $-1/u(k)\approx\eta sk$ for large $k$.

The continuous approximation referred to in Section~\ref{sec:discrete}
replaces $A_k$ and $C_k$ by their means plus Gaussian fluctuations whose
variance scales with $1/b$, and rescales time by $\eta$, which is why it
depends on $(\eta,b)$ only through $r=\eta/b$. Its mean exit time was
obtained by exact quadrature of the one-dimensional first-passage problem.
The derivation of the diffusion and of its exit-time integral is in the
version-5 material and will be included in the companion appendix.

\begin{thebibliography}{99}\small\raggedright
\setlength{\itemsep}{3pt}

\bibitem{cai2026} Y.~Cai. Bypassing minimization bias: a shift-invariant
variance estimator for off-equilibrium local learning coefficients.
Preprint, arXiv:2606.22389, 2026.

\bibitem{cullen2026} B.~Cullen, S.~Estan-Ruiz, R.~Danait and J.~Li. A
basin-selection perspective on grokking via singular learning theory.
Preprint, arXiv:2603.01192 (v3, May 2026).

\bibitem{hennick2026} M.~Hennick and S.~De Baerdemacker. Almost Bayesian: dynamics of SGD through singular learning theory. \emph{ICLR 2026};
arXiv:2503.22478 (v2, March 2026).

\bibitem{hitchcock2025} R.~Hitchcock and J.~Hoogland. From global to local:
a scalable benchmark for local posterior sampling. Preprint,
arXiv:2507.21449 (v2, August 2026).

\bibitem{larsson2026} M.~Larsson. Langevin dynamics along the zero set of
real-analytic potentials. Preprint, arXiv:2608.09840, 2026.

\bibitem{lau2025llc} E.~Lau, Z.~Furman, G.~Wang, D.~Murfet and S.~Wei. The
local learning coefficient: a singularity-aware complexity measure.
\emph{AISTATS 2025}, PMLR 258; arXiv:2308.12108.

\bibitem{mandt2017} S.~Mandt, M.~D.~Hoffman and D.~M.~Blei. Stochastic
gradient descent as approximate Bayesian inference. \emph{Journal of
Machine Learning Research} 18(134):1--35, 2017.

\bibitem{reguera2001} D.~Reguera and J.~M.~Rub\'i. Kinetic equations for
diffusion in the presence of entropic barriers. \emph{Physical Review E}
64:061106, 2001.

\bibitem{watanabe2009} S.~Watanabe. \emph{Algebraic Geometry and
Statistical Learning Theory}. Cambridge University Press, 2009.

\bibitem{zwanzig1992} R.~Zwanzig. Diffusion past an entropy barrier.
\emph{The Journal of Physical Chemistry} 96(10):3926--3930, 1992.

\end{thebibliography}

\end{document}
